
\section{Information dynamics: one theorem, two postulates}
\label{sec:dynamics}

\subsection{Downward jumps are a theorem}

Two distinct statements underpin the downward half of the postulated dynamics.
Conflating them is an easy mistake, and earlier versions of this work made it;
we therefore separate them explicitly.

Recall first the definition of the \emph{mutual information} between the return
vector $X$ and a signal $Y$:
\begin{equation}
\label{eq:mi}
I(X;Y) \;=\; \KL{p_{XY}}{p_X\otimes p_Y}
\;=\; \iint p(x,y)\,\log\frac{p(x,y)}{p(x)p(y)}\,dx\,dy ,
\end{equation}
that is, the divergence of the joint law from the product of the marginals: how
far the joint law is from independence. It is non-negative, symmetric, zero if
and only if $X\perp Y$, and, unlike $h$, invariant under invertible
reparametrisation. For a scalar Gaussian pair with correlation $\rho$,
$I=-\tfrac12\log(1-\rho^2)$.

\begin{proposition}[Fall in average conditional entropy]
\label{prop:conditioning}
Let $X$ be the return vector and $Y$ a signal. Then
\begin{equation*}
\mathbb{E}_Y\!\left[h(X\mid Y)\right] \;=\; h(X) - I(X;Y) \;\le\; h(X),
\end{equation*}
with equality if and only if $X\perp Y$.
\end{proposition}

\begin{proposition}[Monotonicity of the ceiling under constraint]
\label{prop:constraint}
Let $\mathcal{C}\subseteq\mathcal{C}'$ be two constraint sets and
$\mathcal{Q}(\cdot)$ the set of laws that satisfy them. Then
$\mathcal{Q}(\mathcal{C}')\subseteq\mathcal{Q}(\mathcal{C})$ and hence
\begin{equation*}
\sup_{p\,\in\,\mathcal{Q}(\mathcal{C}')} h(p) \;\le\; \sup_{p\,\in\,\mathcal{Q}(\mathcal{C})} h(p).
\end{equation*}
\end{proposition}

Both results are classical \citep{cover2006} and they are not equivalent; the
detailed proof and a worked Gaussian example are in Appendix~\ref{app:mi}. Their
scopes differ along two axes.

Proposition~\ref{prop:conditioning} \emph{quantifies} the fall (exactly
$I(X;Y)$ nats) but only \emph{on average} over realisations of the signal. For a
particular $y$, $h(X\mid Y=y)$ can exceed $h(X)$: a sufficiently confusing signal
leaves the market more uncertain than before. One should therefore expect a jump
component that is mainly, not exclusively, one-sided, which is what
Section~\ref{sec:program} reports.

Proposition~\ref{prop:constraint}, by contrast, holds \emph{for every
realisation}, but it concerns the maximum-entropy \emph{ceiling}, not the actual
entropy, and it quantifies nothing: it says that the ceiling cannot rise, not by
how much it falls.

What the two say together, and what motivates everything that follows: the fall
is \emph{discontinuous}, because the signal arrives at a point in time, and its
expected size in nats is the information content of the news, an amplitude
measured on a scale that depends on neither the asset, the currency nor the
units.

\paragraph{An honesty caveat.} Neither proposition is used to compute anything in
the empirical part. We measure $\JJ$, not $h(X)-\mathbb{E}_Y[h(X\mid Y)]$:
estimating a differential entropy in dimension~20 on a 252-day window is out of
reach (Appendix~\ref{app:units}). The link between the theorem and the
statistics reported later is a line of reasoning, not an estimated identity, and
the reader should treat it as such.

\subsection{Two postulates}

\begin{postulate}[Relaxation]
\label{p:relax}
In the absence of new information, the structural index decays towards a
long-run level at a constant rate:
$d\JJ_t = -\kappa\left(\JJ_t-\bar\JJ\right)dt + \sigma_\JJ\,dW_t$ with $\kappa>0$.
\end{postulate}

The economics behind P\ref{p:relax} is that a constraint goes stale. The news is
priced in, disagreement re-forms, positions disperse, and the cross-section
recovers the freedom the shock took away. Nothing in information theory delivers
this; it is the substantive assumption of the paper, and $\kappa$ is the object
to be estimated.

\paragraph{Why an Ornstein--Uhlenbeck process and not something else?} The choice
is not one of convenience: it is \emph{forced} once three minimal requirements
are imposed, and it is the only process that meets all three.

\begin{enumerate}
\item \emph{Markov.} The future of $\JJ$ depends on the past only through its
current level. This is the parsimony requirement: postulate no memory beyond
what the current state already holds.
\item \emph{Stationary.} $\JJ$ has a finite long-run level and variance. Without
this the index drifts, it is not a state variable, and the regime comparisons of
Section~\ref{sec:program} are meaningless.
\item \emph{Gaussian noise.} Between two information arrivals we claim to know
nothing about the noise beyond its variance: this is the maximum-entropy
principle applied to the noise itself, exactly as it was applied to returns in
Section~\ref{sec:bridge}.
\end{enumerate}

\citet{doob1942}'s theorem states that \textbf{the Ornstein--Uhlenbeck process is
the only continuous-time process that is at once Markov, Gaussian and
stationary}. The three requirements therefore leave no latitude on the diffusive
part: it is determined. The jump term is added precisely to \emph{break}
Gaussianity where the framework says it fails, namely at information arrivals.
The structure of \eqref{eq:ou} is thus an imposed skeleton plus a deliberately
non-Gaussian component, which is exactly the class introduced by
\citet{bns2001}.

\paragraph{What remains arbitrary.} Three things, which must be stated rather
than presented as if they followed from the theorem. The rate $\kappa$ is
assumed \emph{constant}: nothing rules out a crisis being absorbed faster, or
slower, than ordinary news, and the data used here cannot tell. The pull-back is
\emph{linear} in the deviation from the long-run level, whereas a pull-back that
accelerates far from the mean would be just as plausible. Finally, $\JJ$ is
non-negative by construction, whereas an Ornstein--Uhlenbeck process with a
Brownian term can turn negative: the \citet{bns2001} specification avoids this
by dropping the Brownian term and driving the process with a subordinator with
positive jumps, which guarantees a stationary law supported on $\mathbb{R}^+$.
Empirically the issue does not arise ($\JJ$ stays between $2.17$ and
$21.13$~nats on the main panel and between $0.42$ and $1.69$ on the replication
panel, far from the boundary), but formally one should either drop
$\sigma_\JJ$ or write the process on $\log\JJ$.

Nor is the postulate innocent as a matter of measurement: an estimator with
memory will manufacture decay even with no economic relaxation at all, which is
why Section~\ref{sec:program} measures the decay that a memoryless world already
produces.

\begin{postulate}[Asymmetry of the constraint]
\label{p:sign}
The constraint imposed by information is one-sided in the direction of losses,
so that an entropy jump comes with negative skewness of the market factor.
\end{postulate}

\paragraph{Why this postulate is logically necessary.} The index $\JJ$ is an
\emph{even} function of skewness: flipping the sign of skewness leaves its value
unchanged. The reason is immediate. The shape channel is
$\Dun(z)=\KL{p_z}{\varphi}$, and the normal law $\varphi$ is symmetric; so
replacing $z$ by $-z$ leaves the reference law unchanged, and the divergence too:
\begin{equation}
\label{eq:even}
\Dun(-z)=\Dun(z),\qquad\text{and therefore}\qquad \JJ(-x)=\JJ(x).
\end{equation}
The Edgeworth expansion \eqref{eq:edgeworth} makes it visible at a glance:
skewness $S$ enters squared.

The consequence is structural, not anecdotal. \textbf{A distribution with a long
left tail (a crash) and its mirror image with a long right tail (a melt-up) have
exactly the same $\JJ$.} The index measures the \emph{size} of the information,
never its direction. No refinement of the estimator will change this: it is a
property of the object measured. If the paper is to say anything about direction,
it must therefore add an external ingredient. P\ref{p:sign} is that ingredient,
and it is the smallest possible: a single statement, about the sign, added to a
framework that by construction produces none.

\paragraph{What it rests on.} Not on a mechanism we would assert, but on an
established empirical regularity: the \emph{leverage effect} in the sense of
\citet{black1976}, namely that negative returns are followed by larger rises in
volatility than positive returns of the same size. It is one of the most robust
stylised facts in empirical finance \citep{christie1982,engle1993,nelson1991}.
Translated into the language of this paper: for an equal quantity of
information, the market's response is larger on the way down than on the way up,
so a jump in $\JJ$ comes with negative market skewness.

We deliberately abstain from choosing among the competing causal explanations of
this effect (financial leverage proper, volatility feedback, or asymmetric
withdrawal of liquidity on the sell side). The postulate depends on none of them:
it asserts only the asymmetry, not its cause. It is separately testable (H5) and
separately refutable, and Section~\ref{sec:program} concludes that it is not
established.

\subsection{The resulting dynamics}

Proposition~\ref{prop:conditioning} together with P\ref{p:relax} gives a
non-Gaussian Ornstein--Uhlenbeck process with one-sided jumps, of the class
introduced by \citet{bns2001} for stochastic volatility:
\begin{equation}
\label{eq:ou}
d\JJ_t \;=\; -\kappa\left(\JJ_t-\bar\JJ\right)dt \;+\; \sigma_\JJ\,dW_t \;+\; \int_{z>0} z\,\mathcal{P}(dt,dz).
\end{equation}
The integral term needs careful reading, because its notation is compact.

\paragraph{What $\mathcal{P}$ is not.} It is neither a Gaussian law nor a
function. The letter $\mathcal N$ is reserved in this document for the normal
law, so the jump measure carries a distinct symbol.

\paragraph{What $\mathcal{P}$ is.} A \emph{random counting measure}, known as a
Poisson random measure. One can picture it as a random scatter of points in the
(date, size) plane: each point is an information event, its abscissa is the date
it occurs and its ordinate the size of the jump it causes. Formally,
$\mathcal{P}(dt,dz)$ equals $1$ if there is a jump in the time interval
$[t,t+dt)$ whose size falls in $[z,z+dz)$, and $0$ otherwise. The integral
$\int_{z>0} z\,\mathcal{P}(dt,dz)$ then reads simply: ``if there is a jump at
date $t$, add its size $z$ to $\JJ$; otherwise add nothing''. It is a compact
way of writing $\sum_{j} z_j$, the sum of the jumps that occurred.

\paragraph{What characterises it.} Two ingredients, combined in its
\emph{intensity measure} $\nu(dz)\,dt$, the expected number of jumps of size
between $z$ and $z+dz$ during $dt$: an \emph{arrival rate}
$\lambda=\nu(\mathbb{R}^+)$, in jumps per unit time, and a \emph{size
distribution} $\nu(dz)/\lambda$. On the main panel, v1 estimates
$\lambda\approx5$ jumps per year, a mean size of $0.62$~nat and a maximum of
$3.21$~nats.

\paragraph{What $z$ is.} The size of an individual jump, in nats. Conceptually, it
is the mutual information of the news that caused it, in the sense of
Proposition~\ref{prop:conditioning}. Operationally, it is the change in $\JJ$
picked up by the four-sigma filter of Section~\ref{sec:program}. The two coincide
only approximately, and the paper does not claim to identify them.

\paragraph{Why the integral runs over $z>0$ only.} Because $\JJ$ measures an
entropy \emph{deficit} and, by Proposition~\ref{prop:conditioning}, information
can only increase that deficit in expectation. Within this framework nothing
creates a \emph{sudden increase} in entropy: entropy is restored only by the slow
relaxation of the first term, never by a jump. Hence one-sided jumps.

Two caveats, however. Proposition~\ref{prop:conditioning} is a statement
\emph{on average}: a particularly confusing realisation can raise entropy, so a
faithful specification would allow a small probability of negative jumps rather
than forbid them. And the Poisson measure assumes independent arrivals, whereas
the event study of v1 shows clear clustering of jumps; a self-exciting
Hawkes-type process would be more faithful.

Note finally that the diffusion coefficient is written $\sigma_\JJ$ rather than
$\eta$, the latter being reserved for the radius of stress scenarios
(Section~\ref{sec:stress}). Read on entropy itself rather than on $\JJ$, the sign
flips: entropy falls by jumps and recovers by diffusion. The observable
predictions are sharp, and they are what Section~\ref{sec:program} tests: the
distribution of $\Delta\JJ$ is right-skewed; large moves are up and small moves
down; and after a positive jump the index decays with half-life $\ln 2/\kappa$.

\subsection{Why fat tails, correlation and jumps arrive together}

The framework yields a prediction that no volatility or correlation model
produces on its own. It follows in four steps, without technical vocabulary.

\begin{enumerate}
\item Important news concerns \emph{all} assets at once: the same information
constrains each of them.
\item If the news is bad, all assets therefore fall \emph{on the same day}.
\item That day also appears, taken in isolation, in each asset's own history: for
that asset it is an unusually large loss.
\item Accumulating such days distorts each asset's return distribution: the loss
side lengthens and thickens more than the gain side.
\end{enumerate}

Step 2 is what the dependence channel measures: assets move together, measured
correlation rises, $-H^{dep}$ rises. Step 4 is what the shape channel measures:
extreme losses more numerous and larger than extreme gains, that is higher
kurtosis and negative skewness, so $\DD$ rises. The two channels rise together
because they describe the same day seen from two places: the cross-section for
one, each asset's distribution for the other.

In calm markets, notable moves are asset-specific: they do not align the
cross-section and so have no reason to lift both channels at the same time.

The testable form is a \emph{regime-dependent coupling}: the correlation between
$-\Delta H^{dep}$ and $\Delta\DD$ should be close to zero in calm markets and
strongly positive under stress. This is H6, and it is the paper's answer to the
question ``does the link between fat tails, correlation and entropy jumps
require an additional postulate?'': no. Given
Proposition~\ref{prop:conditioning}, the coupling follows from the common nature
of the news. Only the \emph{sign} requires P\ref{p:sign}.
